% ============================================================
% METHOD — Full version for NeurIPS 2026
% ============================================================

\section{SSM Causal Score}
\label{sec:method}

\subsection{Definition}

Given a Mamba model $\mathcal{M}$ trained to minimize the next-step
prediction loss $\mathcal{L} = \frac{1}{T} \sum_t \|\hat{\bX}_{t+1} -
\bX_{t+1}\|^2$ on a $K$-dimensional time series, we define two causal
score extraction methods.

\begin{definition}[Weight-Projection Causal Score]
\label{def:weight_score}
\begin{equation}
    \bS^{\text{W}} = |\bW_{\text{out}} \cdot \bW_{\text{in}}| \in \R^{K \times K},
\end{equation}
where $\bW_{\text{in}} \in \R^{d \times K}$ is the input projection and
$\bW_{\text{out}} \in \R^{K \times d}$ is the output projection of the
trained model, and $|\cdot|$ denotes element-wise absolute value.
\end{definition}

The entry $\bS^{\text{W}}_{i,j}$ quantifies the magnitude of the learned
linear pathway from input variable $j$ to output variable $i$ through the
model's $d$-dimensional hidden space. This score requires only reading two
weight matrices---no forward pass, no gradient computation.

\begin{definition}[Jacobian Causal Score]
\label{def:jacobian_score}
\begin{equation}
    \bS^{\text{J}}_{i,j} = \frac{1}{T} \sum_{t=1}^{T}
    \left| \frac{\partial \hat{X}_i(t)}{\partial X_j(t)} \right|.
    \label{eq:jacobian_score}
\end{equation}
\end{definition}

The Jacobian score captures the \emph{full computational graph}, including
data-dependent $\bB_t$, $\bC_t$ matrices, gating nonlinearities, and
causal convolution. It requires one forward and backward pass through the
trained model.

\paragraph{Relationship between the two scores.}
The weight-projection score captures the \emph{static} information routing
encoded in $\bW_{\text{in}}$ and $\bW_{\text{out}}$, while the Jacobian
score captures both static and \emph{dynamic} routing through the
data-dependent SSM layers. For linear data (VAR), both scores converge to
the same ranking. For nonlinear data, the Jacobian score is more
expressive but also noisier.

\paragraph{Practical choice.} We recommend $\bS^{\text{W}}$ as the default
due to its simplicity and speed ($O(dK)$ computation, no data pass needed).
We use $\bS^{\text{J}}$ when the Jacobian provides strictly better
AUROC, as determined on a validation split. In our experiments,
$\bS^{\text{W}}$ and $\bS^{\text{J}}$ achieve comparable performance on
all benchmarks (Section~\ref{sec:ablation}).

\subsection{Theoretical Analysis}

\begin{theorem}[SSM Causal Score Recovers VAR($p$) Granger Structure]
\label{thm:var_recovery}
Let $\bX_t = \sum_{\tau=1}^{p} \Phi_\tau \bX_{t-\tau} + \epsilon_t$ be
a stable VAR($p$) process with i.i.d.\ Gaussian noise. Let $\mathcal{M}$
be a Mamba model with model dimension $d \geq Kp$ that achieves the
global minimum of the next-step MSE loss. Then:
\begin{equation}
    \bS^{\text{W}}_{i,j} = 0 \iff X_j \text{ does not Granger-cause } X_i.
\end{equation}
\end{theorem}

\begin{proof}[Proof sketch]
At the global optimum, the model implements the Bayes-optimal predictor
$\hat{\bX}_{t+1}^* = \sum_\tau \Phi_\tau \bX_{t+1-\tau}$. With $d \geq Kp$,
the SSM has sufficient capacity to represent the full VAR($p$) impulse
response. The effective input-to-output mapping factors through
$\bW_{\text{out}} \bW_{\text{in}}$, which under genericity of the
internal weights, inherits the sparsity pattern of the coefficient
matrices: zero signal flow from $j$ to $i$ if and only if
$[\Phi_\tau]_{i,j} = 0$ for all $\tau$. The full proof, handling
rank constraints and the factored form, is in
Appendix~\ref{app:proofs}.
\end{proof}

\begin{proposition}[Layer Aggregation]
\label{prop:layer_agg}
For a multi-layer Mamba model, the per-layer effective mapping
$\bS^{(\ell)} = |\bW^{(\ell)}_{\text{out}} \bW^{(\ell)}_{\text{in}}|$
captures causal effects at the timescale determined by layer $\ell$'s
state dynamics (eigenvalues of $\bar{\bA}^{(\ell)}$). Averaging across
layers yields a score robust to the unknown lag distribution.
\end{proposition}

\subsection{Algorithm}

Algorithm~\ref{alg:ssm_causal} summarizes the complete \methodname{}
pipeline. The method is remarkably simple: train a standard Mamba model,
read two weight matrices, multiply them.

\begin{algorithm}[t]
\caption{\methodname{}: Causal Discovery via SSM Internals}
\label{alg:ssm_causal}
\begin{algorithmic}[1]
\REQUIRE Multivariate time series $\bX \in \R^{T \times K}$
\ENSURE Estimated causal graph $\hat{G}$, strength matrix $\bS$
\STATE \textbf{Train} (or load) Mamba model $\mathcal{M}$ on next-step
       prediction: $\min_\theta \frac{1}{T}\sum_t \|\mathcal{M}(\bX_t) -
       \bX_{t+1}\|^2$
\STATE \textbf{Extract} $\bW_{\text{in}} \in \R^{d \times K}$ and
       $\bW_{\text{out}} \in \R^{K \times d}$ from $\mathcal{M}$
\STATE \textbf{Compute} $\bS = |\bW_{\text{out}} \cdot \bW_{\text{in}}|
       \in \R^{K \times K}$
\STATE \textbf{Zero diagonal}: $\bS_{i,i} \leftarrow 0$ for all $i$
       \hfill \COMMENT{Remove self-loops}
\STATE \textbf{Normalize}: $\bS \leftarrow \bS / \max(\bS)$
\STATE \textbf{Threshold}: $\hat{G}_{i,j} = \mathbb{1}[\bS_{i,j} > \tau]$
       \hfill \COMMENT{$\tau$ via Otsu or validation}
\RETURN $\hat{G}$, $\bS$
\end{algorithmic}
\end{algorithm}

\paragraph{Computational complexity.}
Steps 2--6 are $O(dK^2)$---a single matrix multiplication and
element-wise operations. The dominant cost is Step 1 (training),
which is $O(TdNL)$ for a Mamba model with $L$ layers and state
dimension $N$. Crucially, this training is \emph{shared}: once trained
for prediction, the model can be queried for causal structure
\emph{without retraining}. In contrast, cLSTM requires training $K$
separate networks with group-lasso, costing $O(K \cdot Td)$ total.
PCMCI requires $O(K^2 \tau_{\max})$ conditional independence tests,
each involving regression over the full time series.

\paragraph{Normalization and thresholding.}
We normalize $\bS$ by its maximum entry to obtain scores in $[0, 1]$,
ensuring comparability across datasets of different scales. To convert
the continuous score matrix to a binary causal graph, we use Otsu's
method~\citep{otsu1979threshold}, which finds the threshold $\tau$ that
maximizes between-class variance of the score values. This is fully
unsupervised---no knowledge of the true graph is needed. Alternative
thresholds (percentile-based, gap heuristic) yield similar results
(Appendix~\ref{app:additional}). For primary evaluation, we report
threshold-free metrics (AUROC, AUPRC) that are invariant to the
threshold choice.
